\documentclass[10pt,letterpaper]{article}
\usepackage{cornell-notes}

\title{Clause 23.04.03.07.03: Basic String Assign}
\author{}
\date{}
\setDocTitle{Clause 23.04.03.07.03: Basic String Assign}
\setDocSubtitle{C++ 2024}
\setDocOwner{C++ 2024 Cornell Notes}
\setDocDate{\today}
\setCornellCollection{C++ 2024 Cornell Notes}
\setCornellUnitType{Clause}
\setCornellUnitNumber{23.04.03.07.03}
\setCornellUnitTitle{Basic String Assign}
\hypersetup{pdftitle={Clause 23.04.03.07.03: Basic String Assign},pdfsubject={Cornell notes},pdfauthor={}}

\begin{document}
\maketitle
\begin{CornellOverviewBox}{Overview}
\textbf{Classification:} Standard library - Normative\\[2pt]
\textbf{Learning objective:} Understand the rules, terminology, and semantic consequences associated with basic\_string::assign.
\end{CornellOverviewBox}\begin{CornellSummaryBox}{Summary}
{\sffamily\bfseries\color{Accent} Summary}\par\vspace{3pt}Section 23.4.3.7.3 focuses on basic\_string::assign. Apply the specified interface together with its mandates, effects, result conditions, exception behavior, and complexity constraints. When applying it, preserve the distinction between mandatory requirements, recommendations, permissions, and behavior deliberately left undefined, unspecified, or implementation-defined.
\end{CornellSummaryBox}

\end{document}
